\documentclass[conference]{IEEEtran}
\IEEEoverridecommandlockouts
\usepackage{cite}
\usepackage{amsmath,amssymb,amsfonts}
\usepackage{algorithmic}
\usepackage{graphicx}
\usepackage{textcomp}
\usepackage{xcolor}
\usepackage{subcaption}
\usepackage{blindtext}
\def\BibTeX{{\rm B\kern-.05em{\sc i\kern-.025em b}\kern-.08em
    T\kern-.1667em\lower.7ex\hbox{E}\kern-.125emX}}
\begin{document}

\definecolor{twblue}{HTML}{1f497d}
\definecolor{structure}{HTML}{1f497d}
\definecolor{blindtext}{HTML}{D3D3D3}

\title{Comparative Evaluation of SLAM Implementations [TITLE TBD]}
\author{\IEEEauthorblockN{Victor Swekis, [Group Member 2], [Group Member 3], [Group Member 4]}
}

\maketitle


\begin{abstract}
[Context: why SLAM matters for autonomous robot systems.]
[Problem: which SLAM approaches are compared and why.]
[Method: dataset/simulation, SLAM systems, metrics (ATE, RPE).]
[Results: key quantitative findings.]
[Conclusion: interpretation and limitations.]
\end{abstract}

\section{Introduction}
% Motivation, problem statement, contribution of this work.

\section{Methods}
\subsection{SLAM Approaches}
% Theoretical background of each evaluated SLAM implementation
% (at least one per group member).

\subsection{Dataset and Setup}
% Common dataset / simulation used for all SLAM systems.

\subsection{Evaluation Metrics}
% Absolute Trajectory Error (ATE), Relative Pose Error (RPE),
% map overlays (PSNR, IoU for semantic SLAM).

\section{Results}
% Required: overlay plot of the estimated trajectories of all SLAM systems.
% Required: at least one map/reconstruction per SLAM implementation.

\section{Discussion}
% Critical interpretation of system behaviour in the context of theory,
% differences and limitations.

\section{Conclusion}

% \bibliographystyle{IEEEtran}
% \bibliography{refs}

\end{document}